\newpage
\section{第22章：类型重构}

\begin{taplauthorsolutionentrybox}
\phantomsection
\label{sol:author:22.3.9}
\textbf{22.3.9 解答}

主要的算法约束生成规则如下：
\[
\inferrule*[right=CA-Var]
 {x:T\in\Gamma}
 {\Gamma\vdash_Fx:T\mid_F\varnothing}
\qquad
\inferrule*[right=CA-Abs]
 {\Gamma,x:T_1\vdash_Ft_2:T_2\mid_{F'}C}
 {\Gamma\vdash_F\lambda x:T_1.\,t_2:T_1\to T_2\mid_{F'}C}
\]
\[
\inferrule*[right=CA-App]
 {\Gamma\vdash_Ft_1:T_1\mid_{X,F_1}C_1\\
  \Gamma\vdash_{F_1}t_2:T_2\mid_{F_2}C_2}
 {\Gamma\vdash_Ft_1\,t_2:X\mid_{F_2}
  C_1\cup C_2\cup\{T_1=T_2\to X\}}
\]
其余规则类似。

原规则与算法规则的等价性可写成两项。
\begin{enumerate}
  \item \textbf{可靠性：}若
        \(\Gamma\vdash_Ft:T\mid_{F'}C\)，并且 \(\Gamma,t\) 中出现的变量
        均不在 \(F\) 中，则
        \(\Gamma\vdash t:T\mid_{F\setminus F'}C\)。
  \item \textbf{完备性：}若
        \(\Gamma\vdash t:T\mid_XC\)，则存在由 \(X\) 中名称排列成的序列
        \(F\)，使 \(\Gamma\vdash_Ft:T\mid_{\varnothing}C\)。
\end{enumerate}
两项都对推导直接归纳。可靠性的应用情形使用以下事实：若
\(\Gamma,t\) 中出现的类型变量不在 \(F\) 中，且
\(\Gamma\vdash_Ft:T\mid_{F'}C\)，那么 \(T,C\) 中出现的类型变量不在
\(F\setminus F'\) 中。完备性的应用情形使用以下事实：若
\(\Gamma\vdash_Ft:T\mid_{F'}C\)，且 \(G\) 是一列全新的变量名，则
\(\Gamma\vdash_{F,G}t:T\mid_{F',G}C\)。

\taplsolutionbacklink{ex:22.3.9}{返回练习22.3.9}
\end{taplauthorsolutionentrybox}

\begin{taplauthorsolutionentrybox}
\phantomsection
\label{sol:author:22.3.10}
\textbf{22.3.10 解答}

用类型对的列表表示约束集后，约束生成算法就是把
\taplnamedref{sol:author:22.3.9}{解答22.3.9}中的推导规则逐项写成程序：
\begin{taplocamlcode}
let rec recon ctx nextuvar t = match t with
    TmVar(fi,i,_) ->
      let tyT = getTypeFromContext fi ctx i in
      (tyT, nextuvar, [])
  | TmAbs(fi,x,tyT1,t2) ->
      let ctx' = addbinding ctx x (VarBind(tyT1)) in
      let (tyT2,nextuvar2,constr2) = recon ctx' nextuvar t2 in
      (TyArr(tyT1,tyT2), nextuvar2, constr2)
  | TmApp(fi,t1,t2) ->
      let (tyT1,nextuvar1,constr1) = recon ctx nextuvar t1 in
      let (tyT2,nextuvar2,constr2) = recon ctx nextuvar1 t2 in
      let NextUVar(tyX,nextuvar') = nextuvar2() in
      let newconstr = [(tyT1,TyArr(tyT2,TyId(tyX)))] in
      (TyId(tyX), nextuvar',
       List.concat [newconstr; constr1; constr2])
  | TmZero(fi) -> (TyNat, nextuvar, [])
  | TmSucc(fi,t1) ->
      let (tyT1,nextuvar1,constr1) = recon ctx nextuvar t1 in
      (TyNat,nextuvar1,(tyT1,TyNat)::constr1)
  | TmPred(fi,t1) ->
      let (tyT1,nextuvar1,constr1) = recon ctx nextuvar t1 in
      (TyNat,nextuvar1,(tyT1,TyNat)::constr1)
  | TmIsZero(fi,t1) ->
      let (tyT1,nextuvar1,constr1) = recon ctx nextuvar t1 in
      (TyBool,nextuvar1,(tyT1,TyNat)::constr1)
  | TmTrue(fi) -> (TyBool,nextuvar,[])
  | TmFalse(fi) -> (TyBool,nextuvar,[])
  | TmIf(fi,t1,t2,t3) ->
      let (tyT1,nextuvar1,constr1) = recon ctx nextuvar t1 in
      let (tyT2,nextuvar2,constr2) = recon ctx nextuvar1 t2 in
      let (tyT3,nextuvar3,constr3) = recon ctx nextuvar2 t3 in
      let newconstr = [(tyT1,TyBool); (tyT2,tyT3)] in
      (tyT3,nextuvar3,
       List.concat [newconstr; constr1; constr2; constr3])
\end{taplocamlcode}
\taplrustcounterpartflowchunkfirst{ch22-constraints}{1}{30}
\taplrustcounterpartflowchunk{ch22-constraints}{31}{49}
\taplrustcounterpartflowchunk{ch22-constraints}{50}{71}
\taplrustcounterpartflowchunk{ch22-constraints}{72}{95}
\taplrustsupportflowchunkfirst{ch22-inference-support}{1}{24}
\taplrustsupportflowchunk{ch22-inference-support}{25}{48}
\taplrustsupportflowchunk{ch22-inference-support}{49}{60}

\taplsolutionbacklink{ex:22.3.10}{返回练习22.3.10}
\end{taplauthorsolutionentrybox}

\begin{taplauthorsolutionentrybox}
\phantomsection
\label{sol:author:22.3.11}
\textbf{22.3.11 解答}

由\cref{fig:11-12}的 T-Fix 可直接得到：
\[
\inferrule*[right=CT-Fix]
 {\Gamma\vdash t_1:T_1\mid_{X_1}C_1\\
  X\notin X_1\cup\fv(T_1)\cup\fv(C_1)}
 {\Gamma\vdash\taplterm{fix}\ t_1:X\mid_{X_1\cup\{X\}}
  C_1\cup\{T_1=X\to X\}}
\]
它先重构 \(t_1\) 的类型 \(T_1\)，再以新变量 \(X\) 要求
\(T_1=X\to X\)，最终把 \(X\) 作为 \(\taplterm{fix}\ t_1\) 的类型。
\texttt{letrec} 的规则可由此规则和 \texttt{letrec} 的派生形式定义得到。

\taplsolutionbacklink{ex:22.3.11}{返回练习22.3.11}
\end{taplauthorsolutionentrybox}

\begin{taplauthorsolutionentrybox}
\phantomsection
\label{sol:author:22.4.3}
\textbf{22.4.3 解答}

\[
\begin{array}{c@{\qquad}c}
\{X=\tapltype{Nat},Y=X\to X\}
 &[X\mapsto\tapltype{Nat},Y\mapsto\tapltype{Nat}\to\tapltype{Nat}]\\
\{\tapltype{Nat}\to\tapltype{Nat}=X\to Y\}
 &[X\mapsto\tapltype{Nat},Y\mapsto\tapltype{Nat}]\\
\{X\to Y=Y\to Z,Z=U\to W\}
 &[X\mapsto U\to W,Y\mapsto U\to W,Z\mapsto U\to W]\\
\{\tapltype{Nat}=\tapltype{Nat}\to Y\}&\text{不可合一}\\
\{Y=\tapltype{Nat}\to Y\}&\text{不可合一}\\
\varnothing&[\,]
\end{array}
\]

\taplsolutionbacklink{ex:22.4.3}{返回练习22.4.3}
\end{taplauthorsolutionentrybox}

\begin{taplauthorsolutionentrybox}
\phantomsection
\label{sol:author:22.4.6}
\textbf{22.4.6 解答}

可继续用约束列表表示替换，但要求每个数对左侧都是合一变量。先定义单变量替换、
把整个替换作用于类型、把单变量替换作用于约束集，以及出现检查；合一函数随后
就是\cref{def:22.4.4}伪代码的直接实现：
\begin{taplocamlcode}
let substinty tyX tyT tyS =
  let rec f tyS = match tyS with
      TyArr(tyS1,tyS2) -> TyArr(f tyS1,f tyS2)
    | TyNat -> TyNat
    | TyBool -> TyBool
    | TyId(s) -> if s=tyX then tyT else TyId(s)
  in f tyS

let applysubst constr tyT =
  List.fold_left
    (fun tyS (TyId(tyX),tyC2) -> substinty tyX tyC2 tyS)
    tyT (List.rev constr)

let substinconstr tyX tyT constr =
  List.map
    (fun (tyS1,tyS2) ->
      (substinty tyX tyT tyS1, substinty tyX tyT tyS2))
    constr

let occursin tyX tyT =
  let rec o tyT = match tyT with
      TyArr(tyT1,tyT2) -> o tyT1 || o tyT2
    | TyNat | TyBool -> false
    | TyId(s) -> s=tyX
  in o tyT

let unify fi ctx msg constr =
  let rec u constr = match constr with
      [] -> []
    | (tyS,TyId(tyX))::rest ->
        if tyS=TyId(tyX) then u rest
        else if occursin tyX tyS then error fi (msg ^ ": circular constraints")
        else List.append (u (substinconstr tyX tyS rest))
                         [(TyId(tyX),tyS)]
    | (TyId(tyX),tyT)::rest ->
        if tyT=TyId(tyX) then u rest
        else if occursin tyX tyT then error fi (msg ^ ": circular constraints")
        else List.append (u (substinconstr tyX tyT rest))
                         [(TyId(tyX),tyT)]
    | (TyNat,TyNat)::rest -> u rest
    | (TyBool,TyBool)::rest -> u rest
    | (TyArr(tyS1,tyS2),TyArr(tyT1,tyT2))::rest ->
        u ((tyS1,tyT1)::(tyS2,tyT2)::rest)
    | _ -> error fi "Unsolvable constraints"
  in u constr
\end{taplocamlcode}
\taplrustcounterpart{ch22-unify}

这个教学版合一器没有投入太多工作改善错误信息。实际编译器中，清楚解释类型
错误往往是实现类型重构语言时最困难的工程问题之一；可参见 Wand（1986）。

\taplsolutionbacklink{ex:22.4.6}{返回练习22.4.6}
\end{taplauthorsolutionentrybox}

\begin{taplauthorsolutionentrybox}
\phantomsection
\label{sol:author:22.5.6}
\textbf{22.5.6 解答}

把类型重构直接扩展到记录并不容易。主要困难是投影应生成什么约束。朴素规则
会写成
\[
\inferrule*[right=CT-Proj]
 {\Gamma\vdash t:T\mid_X C\\
  Y\notin X\cup\fv(T)\cup\fv(C)}
 {\Gamma\vdash t.l_i:Y\mid_{X\cup\{Y\}}
  C\cup\{T=\{l_i:Y\}\}}
\]
这条规则要求 \(t.l_i\) 的接收者恰好具有单字段记录类型
\(\{l_i:Y\}\)，因而错误地排除含其他字段的记录。

Wand（1987）提出并由 Rémy 等人进一步发展的办法，是增加一类
\emph{行变量}。它的取值不是单个类型，而是整行“字段标签及相应类型”。令
\(\oplus\) 合并字段互不相交的两行，则投影可生成约束
\[
\inferrule*[right=CT-Proj]
 {\Gamma\vdash t:T\mid_XC\qquad
  \rho,\sigma,X\text{ 为新变量}}
 {\Gamma\vdash t.l_i:X\mid_{X\cup\{\rho,\sigma,X\}}
  C\cup\{T=\{\rho\},\rho=(l_i:X)\oplus\sigma\}}
\]
含义是：只要 \(t\) 的记录行 \(\rho\) 包含 \(l_i:X\)，其余字段由 \(\sigma\)
表示，投影结果就具有类型 \(X\)。新的约束含有满足结合律和交换律的
\(\oplus\)，已经不再是普通类型等式；求解时需要一种等式合一算法。

\taplsolutionbacklink{ex:22.5.6}{返回练习22.5.6}
\end{taplauthorsolutionentrybox}
